\subsection{The dimension-two cardinal frontier}

\begin{claim}{A.20}{Proved}{The continuum obstruction and an exact family of thresholds}
Put $\mathfrak c=2^{\aleph_0}$. For $W=S^2$ and every cardinal
$\mu\geq\mathfrak c$, there is no group $G$ with an actual
two-dimensional classifying CW complex and an R2 resolution
$BG\to W$ with bound $\mu$.

More generally, let $W$ be a simply connected two-dimensional CW
complex with nonzero $H_2(W;\mathbb Z)$ of rank $\kappa$, and put
$\nu=\max(\aleph_0,\kappa)$. No such dimension-two R2 resolution
exists for $\mu\geq\nu^{\aleph_0}$.
For every infinite cardinal $\lambda$ satisfying
$\lambda^{\aleph_0}=\lambda$, the particular target
\[
 W_\lambda=\bigvee_{\lambda} S^2
\]
admits a dimension-two R2 resolution with bound $\mu$ exactly
for the cardinals $\mu<\lambda$. All cardinals and cell sets
here belong to the fixed universe $\U$.
\end{claim}
\begin{proof}
An R2 map $t:BG\to W$ is an integral homology isomorphism, by
taking the constant local coefficient system $\mathbb Z$ in its
definition. Since $\pi_1W=1$, its perfect kernel is the whole
group $G$. Therefore
\[
 H_1(G;\mathbb Z)=0,\qquad
 H_2(G;\mathbb Z)\cong H_2(W;\mathbb Z).
 \tag{A.20.1}
\]
The latter group is nonzero, so $G$ is nontrivial: the trivial
group has the point as a classifying space and zero second
homology. Choose $g\ne1$ in $G$. Apply A.19 to its specified
two-dimensional model and to the cardinal $\nu$. It supplies an
actual homomorphism
\[
 f_g:G\longrightarrow U_g,\qquad
 |U_g|\leq\nu^{\aleph_0},\qquad f_g(g)\ne1.
 \tag{A.20.2}
\]
For $\mu\geq\nu^{\aleph_0}$, this is a homomorphism to a group
of size at most $\mu$ that does not vanish on the R2 kernel
$P=G$. This directly contradicts the defining no-small-images
property. For $W=S^2$, the rank in (A.20.1) is one, so
$\nu^{\aleph_0}=\aleph_0^{\aleph_0}=\mathfrak c$.

For $W_\lambda$, the same argument and
$\lambda^{\aleph_0}=\lambda$ exclude every $\mu\geq\lambda$.
For every $\mu<\lambda$, A.13 constructs a simple group
$S_\lambda$ of cardinality $\lambda$, an actual two-dimensional
classifying model, and an acyclic map
\[
 BS_\lambda\longrightarrow W_\lambda
\]
inducing the specified isomorphism on $H_2$. Its kernel is
$S_\lambda$ itself. For a homomorphism
$S_\lambda\to Q$ with $|Q|\leq\mu<\lambda$, simplicity makes
the kernel either the whole group or the trivial group. A
trivial kernel would give an injection of a set of cardinality
$\lambda$ into one of smaller cardinality. Thus the homomorphism
is trivial. This verifies every R2 requirement and proves the
asserted exact threshold for this family.
\end{proof}

\paragraph{State of the remaining question.}
A.20 refutes unrestricted dimension preservation in dimension
two, even for the fixed target $S^2$. A.10 proves existence for
every finite bound. For this fixed target the interval
\[
 \aleph_0\leq\mu<\mathfrak c
 \tag{A.20.3}
\]
has not been settled here. In particular the case
$\mu=\aleph_0$, allowing an uncountable group $G$, remains
undecided. The theorem for the family $W_\lambda$ proves an
exact cardinal relation for those specified targets; it does
not fill (A.20.3). The residual-cardinality proof A.19 uses the
finite nonsingular-system theorem of
\citet{D_OpenAI2026Nonsingular}, an unrefereed preprint whose
complete proof was checked as recorded in the source audit.
